\subsubsection{Modalidad regional, separación y control de réplica}
\label{sec:sd4-dr-modalidad-replica}

\textbf{DR-4 1.0: autoridad y compromiso seleccionados.} La recuperación regional será activo-pasiva, con Producción en São Paulo (\texttt{sa-east-1}) y pasivo preparado en México Central (\texttt{mx-central-1}). Multi-AZ resuelve una falla zonal dentro de cada región, no crea el segundo sitio. La modalidad \emph{warm standby} mantiene las dieciséis tareas mínimas, dos almacenes RDS replicados, objetos, red, claves y controles operativos; el pasivo procesa comprobaciones y recepción de réplicas sin aceptar escrituras de negocio ni ejecutar avisos externos. El compromiso regional es RTO máximo de cuatro horas y RPO máximo de quince minutos para los servicios críticos. El RTO se cuenta desde la interrupción del servicio hasta la restitución funcional, e incluye detección, decisión, aislamiento, promoción, capacidad, identidad y validación; no empieza sólo al autorizar el cambio DNS. Las veinticuatro horas de autonomía local y los treinta minutos de sincronización de C07 pertenecen a otra cadena y mantienen su exigencia. Estos objetivos requieren aceptación futura; esta sección no afirma tiempos ya medidos \parencites[RT-07.01--04, p.~17]{pucv-transversal}[cap.~15]{caso07}.

\textbf{Costo y complejidad operacional.} El pasivo conserva 18 vCPU y 72 GiB de tareas. Producción normal utiliza exactamente esa misma capacidad: no se declara ahorro de cómputo frente a un duplicado normal. Frente a duplicar permanentemente el peak de 32 vCPU/128 GiB, la reserva mínima evita 14 vCPU/56 GiB, un 43,75\,\%; frente a 52 vCPU/208 GiB de crecimiento, evita 34 vCPU/136 GiB, un 65,38\,\%. Son diferencias de capacidad permanente, no descuentos de la factura: el pasivo se amplía al perfil productivo vigente antes de recibir tráfico, y existen consumo de pruebas, almacenamiento, Multi-AZ de las dos réplicas, red, transferencia, monitoreo e identidad. RDS conserva clase/volumen del origen según el plan de capacidad, sin un ahorro ficticio de base de datos. OE valoriza esos recursos por región y horas; esta justificación no introduce precios en el SD. En particular, la reserva adicional MASS/ENV de 170/510 GiB también tiene copia regional y versiones, distintas del subtotal de origen.

Activo-activo exigiría resolver escritura concurrente, asignación exclusiva de amarras/puestos, efectos contables y externos y revocaciones bajo partición; ni Multi-AZ ni dos pools de identidad suministran ese arbitraje. La autoridad única conserva los contratos transaccionales e idempotentes de SD5 y limita esa complejidad. A cambio, activo-pasivo requiere promoción dirigida, capacidad preparada, recuperación de identidad y fencing efectivo; no ofrece continuidad regional sin interrupción ni conmutación automática. Recuperar sólo desde respaldo reduce capacidad permanente, pero incorpora provisión y restauración de todo el volumen y no constituye la modalidad seleccionada. La guía completa de conmutación, aislamiento y retorno se conserva como desarrollo de RT-07.05/06; elegir una modalidad no acredita esos requisitos por sí solo.

\textbf{Distancia declarada y alcance del método.}
\label{sec:sd4-dr-distancia}
Se declara una separación geográfica orientativa de 7.614 km entre las referencias urbanas de São Paulo y Santiago de Querétaro. Se usan las coordenadas GeoNames de São Paulo, ID 3448439, $(-23,54750^\circ,-46,63611^\circ)$, y Santiago de Querétaro, ID 3991164, $(20,58806^\circ,-100,38806^\circ)$, con radio terrestre medio $R=6.371$ km y distancia de círculo máximo:
\[
 d=2R\arcsin\sqrt{\sin^2((\phi_2-\phi_1)/2)+\cos\phi_1\cos\phi_2\sin^2((\lambda_2-\lambda_1)/2)}.
\]
Las coordenadas son referencias públicas de ciudades, no direcciones o coordenadas de centros de datos AWS. São Paulo denomina la región primaria y la comunicación oficial del lanzamiento de México vincula la infraestructura regional con Querétaro; el proveedor no publica en esas fuentes las posiciones de cada zona. El resultado es 7.613,84 km antes del redondeo y no representa distancia de cable ni latencia. La separación intercontinental impide que un incendio, inundación local, corte de distribución eléctrica o evacuación de un único recinto afecte físicamente ambos sitios por ese mismo evento. No se deriva independencia absoluta de fallas por multiplicar kilómetros \parencites{l14-geonames-sp,l14-geonames-qro,l14-aws-mexico}.

\textbf{Amenazas comunes y controles de recuperación.} La separación regional limita fenómenos físicos locales, pero conserva dependencias comunes que se deben evaluar conforme a RT-07.02.

Un incendio, inundación, sismo o corte de energía local queda separado geográficamente; la distancia no certifica resistencia de las instalaciones brasileñas o mexicanas. Un fenómeno amplio puede afectar rutas y logística. Se utilizan dos zonas por región y se comprobará el alcance de continuidad del proveedor mediante retirada de una región completa en el ensayo, incluyendo la conectividad con el recinto.

Ambas regiones comparten AWS y políticas de organización, con riesgo de fallo global o suspensión de cuenta. Se segregan cuentas productiva/DR y de custodia, permisos de operación regional, doble control de cambios de organización y recuperación bajo custodia independiente. No se promete supervivencia de DR a una indisponibilidad total de AWS.

IAM/organización, Route 53/CloudFront, el proveedor federado y los emisores pueden ser dependencias comunes aunque los pools sean regionales. Las credenciales de emergencia son nominativas y se custodian fuera del pool primario; los certificados se preparan en DR y se prueba desde redes externas con resolución y caché. Una caída común de federación o del canal de recuperación conserva su limitación de identidad registrada.

Borrar una clave, distribuir un secreto inválido o desplegar la misma falla puede inutilizar ambas regiones. Las claves de cifrado regionales, políticas y custodias separadas conservan versiones recuperables y el artefacto anterior. Los cambios se ensayan antes de promover y la promoción se bloquea si falla el descifrado o la correspondencia de manifiestos.

La réplica propaga borrados, corrupción y cambios maliciosos; la distancia no los filtra. Se mantienen copias protegidas y restauración aislada conforme a RT-07.09--14. Un incidente de integridad exige detener réplica o promoción, conservar versiones y reaplicar las restricciones válidas.

Dos operadores pueden compartir ducto, mayorista o salida internacional; ambos enlaces de la marina pueden caer aunque México funcione. La matriz de rutas y fallas se verifica con los operadores; se mantienen túneles a ambas regiones y se ensaya con las dos WAN retiradas. La autonomía local conserva su diseño propio: DR no sustituye cobertura de terreno ni resuelve todos los permisos offline.


\textbf{REG-REP-4: población íntegra de réplica.}
\label{sec:sd4-dr-replicacion}
SRE mantiene un registro versionado por conjunto y destino: propietario, clasificación, servicio de origen/destino, región/cuenta, mecanismo, secuencia o versión confirmada, última observación válida, error, edad del pendiente más antiguo y responsable. Incluye los nueve dominios de SD5, ambos almacenes RDS, evidencia y originales, copias ENV, staging MASS aceptado, hechos de integración, auditoría, permisos, consentimiento, eliminaciones y revocaciones. Una tabla SQL que contiene un puntero no acredita copia del objeto. Tampoco se acepta una cola vacía como prueba de que todas las operaciones durables fueron entregadas.

\begin{enumerate}
\item \textbf{RDS transaccional y proyecciones.} Se usan dos réplicas asíncronas entre regiones. Cada almacén emite una transacción de latido cada treinta segundos con secuencia e instante UTC, sin dato personal, leída desde la réplica cada sesenta segundos. El observador calcula edad desde la última secuencia efectivamente visible y verifica avance; compara también LSN recibido/reproducido y bytes pendientes cuando estén expuestos con privilegio de monitoreo. Se registra \texttt{ReplicaLag} cada minuto como métrica complementaria. En PostgreSQL la edad de la última transacción puede crecer con el origen ocioso; por eso se distingue latido propio de la métrica nativa, y un valor inválido o ausencia de muestra nunca se transforma en atraso cero \parencite{l14-rds-lag}.
\item \textbf{Objetos S3.} Cada bucket/regla tiene versionado y CRR hacia México; se habilitan métricas de réplica y eventos de fallo. La réplica de objetos históricos requiere carga inicial mediante Batch Replication, verificación de versiones y manifiesto; habilitar CRR no copia por sí solo todo lo anterior. Cada minuto se capturan atraso por regla, bytes/operaciones pendientes y fallos; se usan nombres de métricas efectivamente disponibles en la configuración habilitada. Además, el manifiesto propio compara objeto, versión, tamaño, huella y estado de destino para cada objeto crítico nuevo, con un objeto de latido por minuto y lectura de destino. Se registran \texttt{PENDING}, \texttt{COMPLETED}, \texttt{FAILED} y \texttt{REPLICA}; una confirmación de SQL no sustituye \texttt{COMPLETED}. Borrados, retención/Object Lock y claves se configuran y comprueban por regla; una eliminación autorizada tiene tombstone durable y se reaplica antes de servir, sin borrar evidencia retenida por otra obligación. La réplica es asíncrona: una métrica baja o un objetivo comercial de réplica no garantiza que todos los objetos estén siempre dentro de quince minutos \parencite{l14-s3-replica}.
\item \textbf{Eventos, colas y resultados externos.} EventBridge, SQS y SNS son transporte regional. DR prepara recursos equivalentes desde IaC, sin afirmar que AWS replique automáticamente sus mensajes. Cada hecho aceptado, orden pendiente, intento, destinatario, recibo, DLQ y resultado incierto se persiste por el propietario con inbox/efecto/outbox en la misma transacción de SD5. El transportador no es la única copia de un hecho crítico. En DR se reconstruyen pendientes a partir del corte SQL promovido y se aplican IDs/huellas, marcas de efecto de toda la retención y versiones del agregado; un mensaje tardío no genera otro cargo, combustible, asignación o envío. Las respuestas inciertas de terceros se consultan por ID antes de reintentar; si el tercero carece de consulta/idempotencia, el efecto queda detenido para conciliación competente, sin inventar un fallo definitivo. Cada minuto se compara secuencia durable de origen/destino por flujo y edad de pendientes; backlog de entrega se separa de lag de réplica. Un webhook sin persistencia confirmada no recibe acuse exitoso. Captura que exista sólo en un bus, cron o memoria obliga a corregir ese productor antes de habilitarlo.
\item \textbf{Borde y equipos incomunicados.} La réplica síncrona del par local, las colas durables de terminales y la exportación por WAN mantienen sus identificadores y acuses. DR observa por separado último hecho recibido y confirmado por el recinto, edad del enlace y cantidad/tamaño no sincronizados conocidos. Un dispositivo sin comunicación tiene estado desconocido; no figura con atraso cero. Sus hechos aún no recibidos por nube no entran artificialmente en el RPO regional: mantienen la obligación C07 de captura sin pérdida, autonomía y posterior sincronización. La reconstrucción DR no purga la cola local ni le asigna una nueva autoridad de amarra.
\item \textbf{Identidad y permisos.} Maestro Persona, relación estable persona-emisor, roles, consentimiento, restricciones y revocaciones se incluyen en SQL y en la comparación de secuencias. El pool Cognito de contingencia tiene configuración propia; se verifican versiones de cliente, claims, políticas, federación y callback. Cognito conserva perfiles en la región del pool; no se afirma copia nativa de contraseñas ni sesiones hacia México. Las personas federadas necesitan disponibilidad del IdP; las demás requieren recuperación/enrolamiento validado con canal verificado. El cambio de emisor invalida sesiones anteriores y exige reautenticación, con permisos/revocaciones del corte promovido. Falta desarrollar la recuperación íntegra del estado exclusivo del pool, MFA y sus canales para todos los perfiles, con medición de atraso o configuración invariable verificable. Esta brecha impide cerrar RT-07.03: una réplica de roles SQL no replica todo el estado de identidad \parencite{l14-cognito-region}.
\item \textbf{Secretos, claves y artefactos.} Secrets Manager replica los secretos que deben compartir valor y la rotación de origen; endpoints y credenciales regionales distintos se aprovisionan separadamente. Se consulta estado de réplica y versión cada minuto, se comprueba descifrado y se impide promover consumidores con versión no compatible. La política de KMS, alias, grants, habilitación y custodia se verifican por región: incluso las claves multirregión requieren administración independiente; no se asume réplica de política por compartir material. No se selecciona aquí una nueva clave multirregión para sustituir sin migración las claves regionales existentes. IaC, contratos, ECR, SBOM y firmas se distribuyen por versión/digest; un cambio no se habilita en Producción hasta verificar su presencia utilizable en DR y el artefacto de retorno. ECR no replica contenido preexistente por activar una regla: se hace carga inicial y verificación por digest. Se sondea cada cinco minutos el manifiesto liberado y de retorno desde México; son réplicas al cambiar versión, no datos que mutan cada minuto. Credenciales de terceros, certificados, políticas y claves antiguas de objetos retenidos requieren mapa completo de custodia/recuperación, todavía pendiente; un inventario de nombres no acredita valor utilizable \parencites{l14-secrets-replica,l14-kms-region,l14-ecr-replica}.
\end{enumerate}

\textbf{Recursos del observador.} REG-REP-4 se ejecuta como trabajo propio de la unidad Auditoría ya inventariada, con un lease regional durable y un único ejecutor activo por región; las dos tareas permiten relevo sin multiplicar sondeos. No añade una novena unidad, servidor o licencia. Se asignan por tarea hasta 256 MiB y 0,05 vCPU del margen existente de 1 GiB/0,30 vCPU, separados del 60\,\% de CPU de negocio y del 10\,\% de agentes. Cada ciclo de sesenta segundos tiene un presupuesto inicial de 120 llamadas de control de respuesta acotada, con hipótesis de 0,010 s CPU por llamada y hasta un segundo CPU adicional de comparación: $(120\times0,010+1)/60=0,0367$ vCPU, inferior a 0,05. Se comprueba CPU, memoria, latencia de API y ciclo real; una llamada que tarda en red no se equipara a CPU y usa concurrencia máxima de cuatro, timeout cinco segundos y dos reintentos dentro del ciclo. Una fase incompleta conserva error y cursor, no informa atraso cero. Crecimiento asigna 0,15 vCPU y 768 MiB de esos márgenes, sin quitar presupuesto de negocio; si la carga exige más, se amplía Auditoría antes de habilitar nuevos conjuntos, preservando capacidad de zona superviviente.

Las métricas se obtienen en lotes y los manifiestos se recorren incrementalmente con cursores por secuencia/versión; no se hace inventario completo de todos los objetos cada minuto. El seguimiento de cada objeto crítico recién aceptado se conserva en SQL antes del acuse de custodia, y se compara en destino sin descargar contenido para cada sondeo. La reconciliación completa por versiones se programa diariamente y al cambiar reglas, con la cuota compartida de tareas pesadas y prioridad inferior a captura. Los dos almacenes agregan en total cuatro pequeños latidos SQL por minuto; los latidos S3 son de 1 KiB, uno por minuto y región, con siete días de retención, inferiores a 20 MiB de contenido total antes de versiones/metadatos. Manifiestos, errores y muestras usan las reservas de SQL/AUD ya existentes y deben entrar en su dimensionamiento integral, aún abierto: no se inventa capacidad total por medir estos latidos. Métricas, API, eventos y alarmas nativas se valorizan en OE por región y retención bajo observabilidad existente; el canal de guardia usa la operación SRE contractual, sin afirmar una contratación ya ejecutada. La aceptación retira una tarea/zona y comprueba que el lease reanuda, los conjuntos se recorren dentro del minuto y las alertas conservan destinatario/acuse. Superar el presupuesto o saltar objetos exige ampliar/repartir la observación con control de cambios, nunca aumentar silenciosamente el plazo de atraso.

\textbf{Muestreo, umbrales y respuesta de operación.} El observador se ejecuta en cada región sin depender de la API de monitoreo de São Paulo para ver México. Consolida cada minuto los flujos dinámicos anteriores con UTC, incertidumbre de reloj y origen de medición. Si la incertidumbre supera cinco segundos se declara muestra no confiable y se corrige sincronización. En flujos críticos, atraso mayor de 120 s durante dos muestras genera aviso a SRE titular/suplente; mayor de 300 s durante dos muestras genera incidencia al responsable de continuidad y Seguridad; mayor de 600 s en una muestra dispara escalamiento urgente, suspensión de liberaciones y de admisión de nuevos trabajos masivos/analíticos, con evaluación de escritura central antes de llegar a 900 s. Llegar a 900 s, fallar réplica/descifrado, perder secuencia, detectar destino incorrecto o no tener dos muestras consecutivas válidas es incidencia crítica y bloquea promoción automática o declaración de RPO conforme. No se interrumpe ni purga captura local para ocultar atraso; tampoco se presenta una suspensión tardía de escritura como garantía retroactiva de RPO.

La consola DR muestra por conjunto el corte confirmado, pendientes, errores y lag, además de disponibilidad del propio observador. SRE atiende 24\,$\times$\,7; los avisos se transmiten por canales operativos contratados disponibles desde México y se conservan en la consola aunque falle un proveedor de mensajes. Un fallo de recepción por el titular escala al suplente a los cinco minutos y al responsable de continuidad a los diez; éste informa al contacto operacional del CLIENTE y decide contingencia, sin exigirle administrar AWS. Se verifican permisos, cuotas, KMS, conectividad, espacio y errores de origen/destino; reintento o Batch Replication se dirige por manifiesto, preservando versiones e integridad. El retorno a estado sano exige tres muestras válidas con atraso menor de 120 s y sin fallos abiertos, comprobación del corte común y cierre de incidencia; no sólo borrar una alarma. SRE conserva muestras crudas, alertas, acuses y acciones en AUD-5, bajo la retención aplicable. El canal de guardia y las métricas que se consultan deben inventariarse en T-11/OE; esta sección no declara una contratación ejecutada.

\textbf{Aceptación y límite del desarrollo.} Se comprobará por flujo pérdida de transporte, objeto fallido, reloj incorrecto, observador ausente, revocación, rotación interrumpida y efecto externo incierto. La prueba compara secuencias y objetos en ambos destinos, verifica cada destinatario de alarma y reconstruye colas desde el registro durable sin efectos duplicados. Se retira una región y se registra el corte común de datos, objetos y restricciones antes de habilitar tráfico. RTO/RPO de todos los servicios críticos, identidad no federada/MFA, custodia íntegra de secretos/credenciales externas y aislamiento de escrituras bajo partición siguen requiriendo desarrollo coordinado; no se obtiene cumplimiento integral de RT-07.03/04/05 por observar sólo \texttt{ReplicaLag}. Las pruebas futuras son evidencia de aceptación de los controles desarrollados, separada de esas brechas de diseño.
